\section*{ID8469: Nd Natural Azimuthal Nodes: π(χ\_T - 1)}
\paragraph{Metadata.}
PiID: 5995744906617 | RTEID: RTE:P32887.0078:T3.9222 | UUID: f914efe7-2248-5d69-864b-ce41da99d67f

\paragraph{Canonical Statement.}
metry. \#\#\#\#\#\# The Trawin Handshake Constant ( \$\textbackslash{}chi\_T\$ ) The ratio between our engineered segments and natural azimuthal nodes is defined as the **Trawin Handshake Constant** ( \$\textbackslash{}chi\_T \textbackslash{}approx 1.185\$ ). This constant is precisely \$32/27\$ , the Pythagorean minor third. | Metric | Natural Nodes | Engineered Nodes | Significance | | ------ | ------ | ------ | ------ | | **Node Count** | 129,600 | 153,600 | \$\textbackslash{}chi\_T \textbackslash{}approx 1.185\$ Ratio | | **Phase Surplus** | \$0\textasciicircum{}\textbackslash{}circ\$ | \$\textbackslash{}approx 33.34\textasciicircum{}\textbackslash{}circ\$ | \$\textbackslash{}pi(\textbackslash{}chi\_T - 1)\$ weighted surplus | | **Function** | Static Alignment | Vacuum Rectification | Impedance match for ZP extraction |

\paragraph{Canonical Equation.}
\[
\pi(\chi_T - 1)
\]

\paragraph{Dictionary.}
Nd Natural Azimuthal Nodes: π(χ\_T - 1) — π(χ\_T - 1)

\paragraph{Encyclopedia.}
Nd Natural Azimuthal Nodes: π(χ\_T - 1) is a symbolic expression, written π(χ\_T - 1). It introduces a symbol or relation used elsewhere in the framework. It is built from π, T. Within the Trawin Zero Point registry it serves as one element of the framework's mathematical narrative.
